%
% 22-polynome.tex -- Der euklidische Algorithmus für Polynome
%
% (c) 2026 Prof Dr Andreas Müller
%
\subsection{Der euklidische Algorithmus für Polynome
\label{buch:ratint:euklid:subsection:polynome}}
Der euklidische Algorithmus funktioniert auch für Polynome.
Dies ermöglicht, rationale Funktionen effizient zu kürzen.
Später werden wir den erweiterten euklidischen Algorithmus
auch verwenden, um die quadratfreie Faktorisierung und die 
zugehörige Partialbruchzerlegung einer rationalen Funktion
zu finden und um Integrale von rationalen Funktionen auf 
einfachere Form zu reduzieren.

%
% Polynomdivision
%
\subsubsection{Polynomdivision}
Auch für Polynome gibt es eine Division mit Rest.
Zu zwei Polynomen $a(x)$ und $b(x)$ gibt es Polynome $q(x)$ und $r(x)$,
für die
\[
a(x) = b(x)q(x)+r(x)
\]
gilt.
Ausserdem muss der Grad des Restpolynoms $r(x)$ kleiner sein als
der Grad von $b(x)$: $\deg r(x) < \deg b()$.

\begin{beispiel}
\label{buch:ratint:euklid:bsp:polynomdivision}
Der Divisionsalgorithmus liefert für die Polynome
\[
a(x)
=
x^4+x^3-4x^2+x-6
\qquad\text{und}\qquad
b(x)
=
x^2+5x+6
\]
die folgende Rechnung:
\[
\renewcommand{\arraycolsep}{1.2pt}
\begin{array}{rcrcrcrcrcrcrcrcrcrcr}
x^4&+&  x^3&-& 5x^2&+&  x&-& 6&:& x^2&+& 5x&+& 6 &=& x^2 &-& 4x&+& 9\\
x^4&+& 5x^3&+& 6x^2& &   & &  & &    & &   & &   & &     & &   & &  \\
\cline{1-5}
   &-& 4x^3&-&11x^2&+&  x& &  & &    & &   & &   & &     & &   & &  \raisebox{4pt}{\mathstrut}\\
   &-& 4x^3&-&20x^2&-&24x& &  & &    & &   & &   & &     & &   & &  \\
\cline{2-7}
   & &     & & 9x^2&+&25x&-& 6& &    & &   & &   & &     & &   & &  \raisebox{4pt}{\mathstrut}\\
   & &     & & 9x^2&+&45x&+&54& &    & &   & &   & &     & &   & &  \\
\cline{5-9}
   & &     & &     &-&20x&-&60& &    & &   & &   & &     & &   & &  \raisebox{3pt}{\mathstrut}\\
\cline{6-9}
\end{array}
\]
Somit ist der Quotient $q(x)=x^2-4x+9$ und der Rest $r(x)=-20x-60$.
\end{beispiel}

Die Division von Polynomen ist die Basis des Algorithmus zur Bestimmung
des grössten gemeinsamen Teilers.
Es ist aber zu beachten, dass die Division einen Rest liefert,
der nicht unbedingt ein normiertes Polynom ist.
Es kann daher nötig sein, den Rest zu normieren.
Im Beispiel ist der normierte Rest $x+3$.

%
% Gemeinsame Teiler für Polynome
%
\subsubsection{Gemeinsame Teiler für Polynome}
Der euklidische Algorithmus für den grössten gemeinsamen Teiler
funktioniert unverändert auch für Polynome.
Der Polynomgrad des Restes nimmt dabei ständig ab, bis der Rest
0 wird.

\begin{beispiel}
Man finde den grössten gemeinsamen Teiler der Polynome
\[
a(x)
=
x^4+x^3-5x^2+x-6
\qquad\text{und}\qquad
b(x)
=
x^2+5x+6.
\]
Dazu müssen die Quotienten und Reste berechnet werden, der
erste wurde bereits im
Beispiel~\ref{buch:ratint:euklid:bsp:polynomdivision}
berechnet:
\begin{center}
\begin{tabular}{|>{$}r<{$}|>{$}l<{$}>{$}l<{$}|>{$}l<{$}>{$}l<{$}|}
\hline
 k & a_k              & b_k      & q_k                & r_k
\raisebox{2.5pt}{\mathstrut}\\[1pt]
\hline
 0 & x^4+x^3-5x^2+x-6 & x^2+5x+6 & x^2-4x+9           & -20x-60
\raisebox{3pt}{\mathstrut}\\
 1 & x^2+5x+6         & -20x-60  & -\frac{1}{20}(x+2) &       0
\\[2pt]
\hline
\end{tabular}
\end{center}
In der zweiten Zeile lässt sich ablesen, dass der Rest $=0$ ist,
daher ist der grösste gemeinsame Teiler $-20x-60$.
Dieses Polynom ist nicht normiert, da die Koeffizienten rational sind,
kann man durch $-20$ teilen und erhält das normierte Polynom $x+3$.
\end{beispiel}

Das Beispiel zeigt auch, dass der gefundene grösste gemeinsame Teiler
nicht ein normiertes Polynom ist.
Da wir nur an Polynomen mit rationalen Koeffizienten interessiert
sind, kann immer durch den Leitkoeffizient geteilt werden.
Arbeitet man dagegen nur mit ganzzahligen Koeffizienten, wird im 
Beispiel bereits der Schritt $k=2$ nicht mehr durchführbar.

%
% Der erweiterte euklidische Algorithmus für Polynome
%
\subsubsection{Der erweiterte euklidische Algorithmus für Polynome}
Auch der erweiterte euklidische Algorithmus lässt sich übertragen.

\begin{beispiel}
Für die Polynome $a(x) = x^3+2x^2-x-2$ und $b(x)=x^2-x-2$
finde man Polynome $s(x)$ und $t(x)$ mit
$\operatorname{ggT}(a,b)=s(x)a(x)+t(x)b(x)$.
\begin{center}
\renewcommand{\arraycolsep}{1.2pt}
\begin{tabular}{|>{$}r<{$}|>{$}l<{$}>{$}l<{$}|>{$}l<{$}>{$}l<{$}|>{$}l<{$}>{$}l<{$}|}
\hline
  k & a_k(x)        & b_k(x)  & q_k(x)       & r_k(x) & s_k(x) & t_k(x) \raisebox{2.5pt}{\mathstrut}\\[1pt]
\hline
 -1 &               &         &              &        &             1 &                0 \\
 -2 &               &         &              &        &             0 &                1 \\
  0 & x^3+2x^2-2x-2 & x^2-x-2 & x+3          & 4x + 4 &             1 &             -x-3 \\
  1 & x^2-x-2       & 4x+4    & \frac14(x-2) & 0      & -\frac14(x-2) & \frac14(x^2+x-2) \\[2pt]
\hline
\end{tabular}
\end{center}
Tatsächlich können wir nachrechnen, dass
\begin{align*}
s_0(x)
\cdot
a(x)
+
t_0(x)
\cdot
b(x)
&=
1
\cdot
(x^3+2x^2-x-2)
-
(x+3)
\cdot
(x^2-x-2)
\\
&=
(x^3+2x^2-x-2)
-
(x^3+2x^2-5x-6)
\\
&=
4x+4
\intertext{und}
s_1(x)
\cdot
a(x)
+
t_1(x)
\cdot
b(x)
&=
-
\frac14(x-2)
\cdot
(x^3+2x^2-x-2)
+
\frac14(x^2+x-2)
\cdot
(x^2-x-2)
\\
&=
-
\frac14(x^4-5x^2+4)
+
\frac14(x^4-5x^2+4)
\\
&=0.
\end{align*}
Verlangt man ein normiertes Polynom für den grössten gemeinsamen Teiler,
kann man die Koeffizienten in $\mathbb{Q}$ durch $4$ dividieren und 
erhält $\operatorname{ggT}(a,b)=x+1$.
Die Faktoren $s(x)$ und $t(x)$ müssen ebenfalls durch 4 geteilt werden
und es gilt
\[
\frac14
\cdot
(x^3+2x^2-x-2)
-\frac14(x+3)
(x^2-x-2)
=
x+1.
\]

Den grössten gemeinsamen Teiler kann man auch mithilfe der Faktorisierung
\begin{align*}
a(x)
&=
(x-1){\color{darkred}(x+1)}(x+2)
\\
b(x)
&=
(x-2){\color{darkred}(x+1)}
\end{align*}
der Polynome finden.
Der einzige gemeinsame Faktor ist $\operatorname{ggT}(a,b)={\color{darkred}x+1}$.
Aus der Faktorisierung sind jedoch die Faktoren $s(x)$ unt $t(x)$
nicht direkt ablesbar.
\end{beispiel}

